\documentclass[border=10pt]{standalone}
\usepackage{fontspec}
\setmainfont{Noto Serif}[Ligatures=TeXOff]
\XeTeXgenerateactualtext=1
\newfontfamily\bengalifont{Noto Serif Bengali}[Script=Bengali, NFSSFamily=bnfam]
\usepackage[Latin,Bengali]{ucharclasses}
% group-free font switching (as in header-common.tex): robust inside braces, cells and quotes
\makeatletter
\def\BN@saved{\rmdefault}
\def\BN@bn{bnfam}
\newcommand\BNin{\edef\BN@tmp{\f@family}\ifx\BN@tmp\BN@bn\else\let\BN@saved\BN@tmp\fi\fontfamily{bnfam}\selectfont}
\newcommand\BNout{\fontfamily{\BN@saved}\selectfont}
\makeatother
\setTransitionsFor{Bengali}{\BNin}{\BNout}
% Typographic quotes and dashes are not in a ucharclasses class (class 0), so a Bangla word right
% after “ or ‘ would not trigger the Bangla font. Treat them as Latin so that every switch fires.
\XeTeXcharclass"2018=\XeTeXcharclass`A
\XeTeXcharclass"2019=\XeTeXcharclass`A
\XeTeXcharclass"201C=\XeTeXcharclass`A
\XeTeXcharclass"201D=\XeTeXcharclass`A
\XeTeXcharclass"2013=\XeTeXcharclass`A
\XeTeXcharclass"2014=\XeTeXcharclass`A
\XeTeXcharclass"2026=\XeTeXcharclass`A
\usepackage{newunicodechar}
\newunicodechar{।}{{\bengalifont\char"0964}}
\newunicodechar{॥}{{\bengalifont\char"0965}}

\usepackage[table]{xcolor}
\usepackage{booktabs,tabularx,array}
\definecolor{bapilrule}{HTML}{1F2937}
\definecolor{bapilstripe}{HTML}{F3F4F6}
\renewcommand{\arraystretch}{1.28}
\setlength{\tabcolsep}{5pt}

\begin{document}
\begin{minipage}{17cm}
{\large\bfseries সারণি ২২।}\ {\large ডেটা শ্রেণিবিন্যাস ও ব্যবস্থাপনার নিয়ম}\par\vspace{7pt}
\small
\rowcolors{2}{bapilstripe}{white}
\begin{tabularx}{\linewidth}{@{}>{\hsize=0.693\hsize\raggedright\arraybackslash}X>{\hsize=0.990\hsize\raggedright\arraybackslash}X>{\hsize=1.188\hsize\raggedright\arraybackslash}X>{\hsize=1.386\hsize\raggedright\arraybackslash}X>{\hsize=0.743\hsize\raggedright\arraybackslash}X@{}}
\arrayrulecolor{bapilrule}\toprule
\rowcolor{white}\textbf{শ্রেণি} & \textbf{বিবরণ} & \textbf{উদাহরণ} & \textbf{ব্যবস্থাপনার নিয়ম} & \textbf{পাইলটে ব্যবহার} \\
\midrule
\textbf{DC0} উন্মুক্ত & সরকারিভাবে প্রকাশিত ও জনসাধারণের জন্য উদ্দিষ্ট & আইন, গেজেট, বাজেট নথি, পরিকল্পনা, পরিসংখ্যান প্রকাশনা & উৎস-তথ্য ও অখণ্ডতা যাচাই; উৎসের শর্ত সাপেক্ষে পুনর্ব্যবহার & হ্যাঁ \\
\textbf{DC1} অভ্যন্তরীণ & অপ্রকাশিত, ব্যক্তিগত নয় এমন তথ্য & চুক্তির অধীনে বিনিময়কৃত খসড়া প্রতিবেদন, প্রশাসনিক মেটাডেটা & ডেটা-বিনিময় চুক্তি; ভূমিকাভিত্তিক প্রবেশাধিকার; বাংলাদেশে সংরক্ষণ & সীমিত, চুক্তির অধীনে \\
\textbf{DC2} ব্যক্তিগত & ব্যক্তিকে শনাক্তযোগ্য তথ্য & নামসহ পরামর্শ-মতামত, জরিপের মাইক্রোডেটা & ব্যক্তিগত উপাত্ত সুরক্ষা আইন ২০২৬ অনুযায়ী আইনগত ভিত্তি; DPIA; বিশ্লেষণের আগে পরিচয়-অপসারণ; ERRC অনুমোদন & শুধু সম্মতিপ্রাপ্ত ও পরিচয়-অপসারিত \\
\textbf{DC3} সংবেদনশীল বা সংরক্ষিত & সংবেদনশীল ব্যক্তিগত বা নিরাপত্তা-শ্রেণিভুক্ত তথ্য & স্বাস্থ্য, বায়োমেট্রিক, আর্থিক বা শিশুদের ডেটা; নিরাপত্তা-সংক্রান্ত ডেটা & BAPIL-এর কাজে বর্জিত & না \\
\bottomrule
\end{tabularx}
\par\vspace{3pt}{\scriptsize\color{gray} উৎস: BAPIL টেকনিক্যাল ব্লুপ্রিন্ট, সংস্করণ ১.০ (২০২৬)। CC BY 4.0।}
\end{minipage}
\end{document}
